\section{System Architecture and Agentic Workflow}
\label{sec:architecture}

\subsection{Explicit State-Machine Core vs. Frameworks}
\label{sec:statemachine}

A central architectural tenet of our design is the rejection of black-box multi-agent frameworks in favor of an \textbf{explicit Python state machine}. In academic research monitoring, auditability and determinism are paramount: a researcher must understand exactly why a manuscript was delivered, queued, or discarded.

The lifecycle of each candidate paper is tracked by formal states across the system:
\begin{enumerate}
    \item \texttt{ROUTED\_HIGH}: Candidate paper scored at or above $t_{\text{high}}$; automatically scheduled for summarization and digest delivery.
    \item \texttt{ROUTED\_BORDERLINE}: Candidate paper scored between $t_{\text{low}}$ and $t_{\text{high}}$; scheduled for review queue ingestion.
    \item \texttt{DISCARDED}: Candidate paper scored at or below $t_{\text{low}}$; permanently marked as discarded (terminal state, logged with full audit metadata).
    \item \texttt{PENDING\_REVIEW}: Borderline paper persisted in the user's active human review queue.
    \item \texttt{REVIEW\_REJECTED}: Borderline paper explicitly rejected by the human reviewer (terminal state, generates a negative training event).
    \item \texttt{REVIEW\_EXPIRED}: Borderline paper unreviewed after 72 hours (terminal state, retired without generating training labels).
    \item \texttt{DEFERRED\_OVERFLOW}: Borderline paper exceeding the cycle review queue cap (15 items) or high-confidence paper exceeding the digest cap (20 items); logged and terminated without labeling.
    \item \texttt{SUMMARIZE\_PENDING}: Gate-cleared paper awaiting LLM synthesis.
    \item \texttt{DELIVERED}: Synthesized paper published to the user's research digest.
    \item \texttt{REACTED}: User provided interactive feedback (Like $\pmb{+}$ or Dislike $\pmb{-}$) on the delivered digest item.
\end{enumerate}

\begin{figure*}[t]
\centering
\resizebox{0.96\textwidth}{!}{%
\begin{tikzpicture}[
    >=Latex,
    block/.style={rectangle, draw=blue!70!black, fill=blue!8, text width=2.5cm, text centered, rounded corners=3pt, minimum height=1.1cm, font=\small\bfseries},
    decision/.style={diamond, draw=orange!80!black, fill=orange!15, text width=2.1cm, text centered, inner sep=1pt, font=\small\bfseries, aspect=1.3},
    terminal/.style={rectangle, draw=gray!70!black, fill=gray!15, text width=2.3cm, text centered, rounded corners=3pt, minimum height=0.9cm, font=\small},
    line/.style={draw=black!85, -{Latex[length=2.5mm]}, thick},
    lbl/.style={font=\scriptsize, fill=white, inner sep=1.5pt}
]

    % Row 1: Primary Pipeline
    \node [block] (input) {User Profile ($\mathbf{v}_0$)\\ \& arXiv Stream};
    \node [block, right=1.0cm of input] (embed) {Dense Embedding\\($384$-d, $L_2$-norm)};
    \node [decision, right=1.0cm of embed] (gate) {Relevance Gate\\($t_{\text{high}}, t_{\text{low}}$)};

    % Three branches
    \node [block, above right=0.9cm and 1.3cm of gate] (high) {HIGH Route\\Auto-Deliver};
    \node [block, right=1.5cm of gate] (border) {BORDERLINE\\Queue (Cap 15)};
    \node [terminal, below right=0.9cm and 1.3cm of gate] (low) {LOW Route\\Discard \& Audit};

    % Post-gate and Delivery
    \node [block, right=1.3cm of high] (summary) {Summary Agent\\(Interest-Framed)};
    \node [block, right=1.3cm of summary] (digest) {Daily Digest\\(User $\pmb{+}/\pmb{-}$ Rating)};

    % Preference Learning
    \node [block, below=1.5cm of digest] (pref) {Preference Memory\\(\texttt{training\_labels})};
    \node [block, below=0.9cm of pref] (adapt) {Rocchio Centroid $\to$\\LogReg Classifier};

    % Connections
    \path [line] (input) -- (embed);
    \path [line] (embed) -- (gate);

    \path [line] (gate) |- node [lbl, pos=0.7] {$s \ge t_{\text{high}}$} (high.west);
    \path [line] (gate) -- node [lbl] {Ambiguous} (border.west);
    \path [line] (gate) |- node [lbl, pos=0.7] {$s \le t_{\text{low}}$} (low.west);

    \path [line] (high) -- (summary);
    \path [line] (border.east) -- ++(0.6,0) |- node [lbl, pos=0.25] {Accept} (summary.south);
    \path [line] (border.east) -- ++(0.6,0) |- node [lbl, pos=0.75] {Reject / Expire} (pref.west);

    \path [line] (summary) -- (digest);
    \path [line] (digest) -- node [lbl] {Feedback} (pref);
    \path [line] (pref) -- (adapt);

    % Feedback loop back to Gate
    \path [line] (adapt.south) -- ++(0,-0.5) -| node [lbl, pos=0.25] {Updated $\mathbf{v}_{t+1}$ / Thresholds} ([yshift=-0.5cm]low.south) -| (gate.south);

\end{tikzpicture}%
}
\caption{End-to-end operational architecture of the Relevance-Gated Agentic Workflow, illustrating the three-way relevance gate, asynchronous human review escalation, dual-tier preference memory, and interest-framed LLM summarization.}
\label{fig:architecture}
\end{figure*}

\subsection{Three-Tier Runtime Deployment Architecture}
\label{sec:three_tier}

To ensure high performance while maintaining compatibility with lightweight hosting constraints, the system is partitioned into three distinct runtime layers:
\begin{enumerate}
    \item \textbf{Always-On Web Service:} Implemented in FastAPI, Jinja2, and HTMX. Serves the researcher's review queue, research digest, and audit log. Crucially, the web layer \textbf{does not import PyTorch or SentenceTransformers}, maintaining an active memory footprint under 150~MB. LLM summarization for review-accepted papers is handled via lightweight asynchronous background tasks.
    \item \textbf{Scheduled Cycle Worker:} Executes automated monitoring cycles (fetching arXiv feeds, computing embeddings, scoring against active profiles, managing queue caps, updating centroid vectors, and retraining classifiers). Designed to execute headlessly via cron or scheduled background workers.
    \item \textbf{Persistent Data Layer:} Managed PostgreSQL instance with the \texttt{pgvector} extension. Stores document metadata, 384-dimensional dense vectors, auditable routing decisions, and serialized model artifacts. For local development and test environments, an automated fallback engine utilizes SQLite with in-memory NumPy vector operations.
\end{enumerate}

\subsection{Database Schema and Relational Architecture}
\label{sec:schema}

The database schema consists of nine core tables and one materialized SQL view:
\begin{itemize}
    \item \texttt{papers}: Stores canonical arXiv publications, containing \texttt{arxiv\_id} (version-stripped unique primary key), \texttt{title}, \texttt{abstract}, \texttt{submitted\_date}, \texttt{categories} (JSON), \texttt{embedding} (\texttt{vector(384)}), \texttt{embedding\_model}, and \texttt{is\_reference} boolean flag.
    \item \texttt{users}: Pilot researcher records with private \texttt{access\_token} strings for URL-authenticated access.
    \item \texttt{interest\_profiles}: Initial seed interest text and corresponding 384-d seed vector.
    \item \texttt{interest\_vector\_versions}: History of all active and legacy interest vectors for each user, recording the derivation \texttt{method} (\texttt{seed} or \texttt{rocchio}), \texttt{n\_effective\_labels}, threshold parameters ($P_{\text{high}}, P_{\text{low}}, t_{\text{high}}, t_{\text{low}}$), and activation status.
    \item \texttt{cycles}: Execution logs of monitoring runs, tracking \texttt{as\_of\_date}, \texttt{mode} (\texttt{replay} vs \texttt{live}), processing timestamps, and summary statistics.
    \item \texttt{routing\_decisions}: Immutable audit log of every paper evaluated per user. Captures \texttt{cosine\_score}, \texttt{clf\_prob}, \texttt{method} (\texttt{threshold} vs \texttt{classifier}), \texttt{route} (\texttt{HIGH}, \texttt{BORDERLINE}, \texttt{LOW}), \texttt{status}, and decision \texttt{reason}.
    \item \texttt{review\_tasks}: Asynchronous queue entries for borderline papers, tracking presentation timestamps, decisions (\texttt{accept}, \texttt{reject}, \texttt{skip}), and expiration markers.
    \item \texttt{digests} and \texttt{digest\_items}: Curated daily delivery collections, capping delivered papers at 20 per cycle and tracking delivery provenance (\texttt{via='auto'} vs \texttt{via='review'}).
    \item \texttt{feedback\_events}: Append-only ledger of every user reaction (Like $\pmb{+} = 1$, Dislike $\pmb{-} = 0$) and reviewer decision (Accept = 1, Reject = 0).
    \item \texttt{training\_labels} (SQL View): Deterministic resolution of multi-source feedback into one effective training label per (user, paper) pair.
\end{itemize}
